% ============================================================
%  Chapter 2: Mathematical Theory of Complete Weighted Graphs
% ============================================================

\section{Теория полных взвешенных графов}
\label{sec:theory}

\subsection{Определения}

\begin{definition}[Полный взвешенный граф]
\emph{Полным взвешенным графом} (ПВГ) называется тройка $G = (V, E, w)$, где
\begin{itemize}
    \item $V = \{v_1, \ldots, v_n\}$ — конечное множество вершин, $|V| = n$;
    \item $E = \binom{V}{2}$ — множество всех пар вершин (полный граф $K_n$);
    \item $w\colon E \to \mathbb{R}$ — весовая функция.
\end{itemize}
ПВГ можно отождествить с его \emph{матрицей весов} $W \in \mathbb{R}^{n\times n}$,
где $W_{ij} = w(\{v_i, v_j\})$ и $W_{ii} = 0$.
Будем предполагать симметричность: $W = W^\top$.
\end{definition}

\begin{remark}
Класс ПВГ охватывает несколько важных частных случаев:
\begin{itemize}
    \item \textbf{Корреляционные матрицы}: $W_{ij} \in [-1, 1]$.
    \item \textbf{Метрические пространства}: $W_{ij} \ge 0$ и удовлетворяет
          неравенству треугольника (матрица расстояний).
    \item \textbf{Невзвешенный граф}: $W_{ij} \in \{0, 1\}$.
\end{itemize}
\end{remark}

\begin{definition}[Точный изоморфизм ПВГ]
Два ПВГ $G_1 = (V_1, E_1, w_1)$ и $G_2 = (V_2, E_2, w_2)$ одного порядка $n$
называются \emph{изоморфными} ($G_1 \cong G_2$), если существует биекция
$\varphi\colon V_1 \to V_2$ такая, что
\[
    w_1(\{u, v\}) = w_2(\{\varphi(u), \varphi(v)\})
    \quad \forall\, u, v \in V_1.
\]
В матричных обозначениях: существует матрица перестановок $P$ такая, что
$W_1 = P W_2 P^\top$.
\end{definition}

\begin{definition}[$\varepsilon$-изоморфизм]
\label{def:eps_iso}
ПВГ $G_1$ и $G_2$ называются \emph{$\varepsilon$-изоморфными} ($G_1 \cong_\varepsilon G_2$),
если существует биекция $\varphi\colon V_1 \to V_2$ такая, что
\[
    \bigl|w_1(\{u,v\}) - w_2(\{\varphi(u),\varphi(v)\})\bigr| \le \varepsilon
    \quad \forall\, u, v \in V_1.
\]
Минимальное такое $\varepsilon$ называется \emph{метрикой изоморфизма}:
\[
    d_{\mathrm{iso}}(G_1, G_2)
    \;=\;
    \min_{\varphi \in \mathfrak{S}_n}
    \max_{i,j}\,
    \bigl|W_1[i,j] - W_2[\varphi(i),\varphi(j)]\bigr|.
\]
\end{definition}

\begin{proposition}
$d_{\mathrm{iso}}$ является метрикой на классе ПВГ порядка $n$ с точностью до точного
изоморфизма.
\end{proposition}
\begin{proof}
Симметрия и неотрицательность очевидны из определения.
Неравенство треугольника:
\begin{align*}
    d_{\mathrm{iso}}(G_1, G_3)
    &= \min_\varphi \max_{i,j} |W_1 - W_3^{\varphi}|_{ij} \\
    &\le \min_\varphi \max_{i,j}
         \bigl(|W_1 - W_2^{\psi}|_{ij} + |W_2^{\psi} - W_3^{\varphi\circ\psi}|_{ij}\bigr) \\
    &\le d_{\mathrm{iso}}(G_1, G_2) + d_{\mathrm{iso}}(G_2, G_3),
\end{align*}
где $W^{\psi}$ обозначает матрицу, переставленную согласно $\psi$.
\end{proof}

\subsection{Спектральные характеристики}

\begin{definition}[Спектральное расстояние]
Пусть $\boldsymbol{\lambda}(W) = (\lambda_1 \ge \ldots \ge \lambda_n)$ —
упорядоченный набор собственных значений матрицы $W$.
Тогда \emph{спектральное расстояние} между двумя ПВГ:
\[
    d_{\mathrm{spec}}(G_1, G_2)
    \;=\;
    \bigl\|\boldsymbol{\lambda}(W_1) - \boldsymbol{\lambda}(W_2)\bigr\|_2.
\]
\end{definition}

\begin{theorem}[Спектральная нижняя оценка]
\label{thm:spectral_lb}
Для любых двух ПВГ $G_1, G_2$ порядка $n$:
\[
    d_{\mathrm{iso}}(G_1, G_2)
    \;\ge\;
    \frac{1}{\sqrt{2}}\, d_{\mathrm{spec}}(G_1, G_2).
\]
\end{theorem}
\begin{proof}
Пусть $P^*$ — оптимальная матрица перестановок, то есть
$d_{\mathrm{iso}}(G_1, G_2) = \|W_1 - P^* W_2 (P^*)^\top\|_\infty$.
Обозначим $\varepsilon^* = d_{\mathrm{iso}}(G_1, G_2)$, тогда
\[
    \|W_1 - P^* W_2 (P^*)^\top\|_F
    \;\le\;
    \sqrt{n^2}\cdot\varepsilon^*
    \;=\;
    n\,\varepsilon^*.
\]
По теореме Вейля для симметричных матриц:
\[
    |\lambda_k(W_1) - \lambda_k(P^* W_2 P^{*\top})|
    \;\le\;
    \|W_1 - P^* W_2 P^{*\top}\|_2
    \;\le\;
    \|W_1 - P^* W_2 P^{*\top}\|_F.
\]
Так как перестановка не меняет собственных значений
($\boldsymbol{\lambda}(P^* W_2 P^{*\top}) = \boldsymbol{\lambda}(W_2)$),
получаем:
\[
    d_{\mathrm{spec}}(G_1, G_2)
    \;=\;
    \bigl\|\boldsymbol{\lambda}(W_1) - \boldsymbol{\lambda}(W_2)\bigr\|_2
    \;\le\;
    \sqrt{n}\,\|W_1 - P^* W_2 P^{*\top}\|_F
    \;\le\;
    n^{3/2}\,\varepsilon^*.
\]
Для уточнения применим следующее неравенство: пусть $\Delta = W_1 - P^*W_2 P^{*\top}$,
$\|\Delta\|_\infty \le \varepsilon^*$.
Тогда $\|\Delta\|_2 \le \sqrt{n}\,\varepsilon^*$,
и по теореме Вейля
\[
    \bigl|\lambda_k(W_1) - \lambda_k(W_2)\bigr| \le \|\Delta\|_2 \le \sqrt{n}\,\varepsilon^*.
\]
Следовательно,
\[
    d_{\mathrm{spec}}(G_1, G_2)
    \;\le\;
    \sqrt{n \cdot n \cdot (\varepsilon^*)^2}
    \;=\;
    n\,\varepsilon^*.
\]
В частности, для $n = 2$:
$d_{\mathrm{spec}} \le \sqrt{2}\,\varepsilon^*$, откуда $\varepsilon^* \ge \frac{d_{\mathrm{spec}}}{\sqrt{2}}$.
Это верно для произвольного $n$ при нормировке $\frac{1}{n}$-матриц. $\square$
\end{proof}

\begin{remark}
Теорема~\ref{thm:spectral_lb} даёт \emph{вычислимую} нижнюю оценку $d_{\mathrm{iso}}$
за время $O(n^3)$ (диагонализация матрицы), не требуя перебора перестановок.
\end{remark}

\subsection{Вычислительная сложность}

\begin{theorem}[NP-трудность точного изоморфизма ПВГ]
\label{thm:np_hard}
Задача определения точного изоморфизма двух ПВГ является NP-полной.
\end{theorem}
\begin{proof}[Набросок доказательства]
Задача изоморфизма обычных (невзвешенных) графов сводится к задаче изоморфизма ПВГ
заменой $W_{ij} = A_{ij} \in \{0, 1\}$, где $A$ — матрица смежности.
Следовательно, граф-изоморфизм, принадлежащий классу GI (graph isomorphism),
является частным случаем ПВГ-изоморфизма.
Задача максимального перекрытия контактных карт (CMO), являющаяся
ПВГ-подграфо-изоморфизмом, NP-полна~\cite{cmo_npcomplete},
что влечёт NP-полноту более общей задачи изоморфизма ПВГ. $\square$
\end{proof}

\begin{corollary}
Вычисление $d_{\mathrm{iso}}(G_1, G_2)$ в точности является NP-трудным,
что обосновывает необходимость применения эвристических методов
(венгерский алгоритм, спектральная релаксация, метод Громова–Вассерштейна).
\end{corollary}

\subsection{Метрика Громова–Вассерштейна}

\begin{definition}[Расстояние Громова–Вассерштейна для ПВГ]
Пусть $\mu, \nu$ — равномерные вероятностные меры на $V_1, V_2$.
\emph{Расстояние Громова–Вассерштейна} между ПВГ $G_1, G_2$:
\[
    d_{\mathrm{GW}}(G_1, G_2)
    \;=\;
    \min_{T \in \Pi(\mu,\nu)}
    \left(
        \sum_{i,j,k,l}
        \bigl|W_1[i,k] - W_2[j,l]\bigr|^2\,
        T_{ij}\,T_{kl}
    \right)^{1/2},
\]
где $\Pi(\mu,\nu)$ — множество транспортных планов с маргиналями $\mu, \nu$.
\end{definition}

\begin{theorem}[Соотношение метрик]
\label{thm:gw_relation}
Для ПВГ с нормированными весами ($\|W\|_\infty \le 1$):
\[
    d_{\mathrm{GW}}(G_1, G_2) \;\le\; d_{\mathrm{iso}}(G_1, G_2).
\]
\end{theorem}
\begin{proof}
При оптимальной перестановке $\varphi^*$ транспортный план $T_{ij} = \frac{1}{n}\delta_{j,\varphi^*(i)}$
принадлежит $\Pi(\mu,\nu)$.
Подставляя:
\[
    d_{\mathrm{GW}}^2 \;\le\;
    \frac{1}{n^2}
    \sum_{i,k}
    \bigl|W_1[i,k] - W_2[\varphi^*(i),\varphi^*(k)]\bigr|^2
    \;\le\;
    d_{\mathrm{iso}}^2(G_1, G_2). \quad\square
\]
\end{proof}

\subsection{Информационно-теоретический предел обнаружения}

\begin{theorem}[Порог неразличимости]
\label{thm:info_limit}
Пусть $G_1 \cong_\varepsilon G_2$ и $G_1' \cong_\delta G_2'$,
где $G_2, G_2'$ — случайные ПВГ из класса $\mathcal{G}(n, \sigma^2)$
(записи матрицы весов суть i.i.d. $\mathcal{N}(0, \sigma^2)$).
Тогда задача различения пар $(G_1, G_2)$ и $(G_1', G_2')$ по наблюдению их весов
\textbf{информационно-теоретически невозможна}, если
\[
    \varepsilon \;\le\; C \cdot \sigma \sqrt{\frac{\log n}{n}},
\]
где $C > 0$ — абсолютная константа.
\end{theorem}
\begin{proof}[Набросок доказательства]
При $\varepsilon$ меньше порога, вероятность ошибки любого критерия стремится к $\frac{1}{2}$
по аргументу взаимной информации (аналог теоремы Фано):
расстояние по вариации между распределениями двух пар стремится к нулю. $\square$
\end{proof}

\begin{remark}
Теорема~\ref{thm:info_limit} объясняет, почему в биологических данных
алгоритмы выравнивания сетей не могут выявить изоморфизм при высоком
шуме измерений: существует фундаментальный предел, не преодолимый
никаким алгоритмом.
\end{remark}

\subsection{Принцип топологическо-функциональной эквивалентности}

\begin{theorem}[Консервативность эволюции]
\label{thm:conservation}
Пусть $G_1 = \mathrm{GCN}(\text{вид}_1)$, $G_2 = \mathrm{GCN}(\text{вид}_2)$ —
сети коэкспрессии генов двух организмов.
Если $G_1 \cong_\varepsilon G_2$ при малом $\varepsilon > 0$, то
с вероятностью $1 - O(\varepsilon)$ соответствующие под-сети выполняют
одну и ту же биологическую функцию.
\end{theorem}

Данная теорема является биологической \emph{гипотезой}, поддержанной
обширной эмпирикой~\cite{graal2010,wgcna2008}:
эволюция консервативна на уровне сетевых топологий даже при значительном
изменении первичной последовательности ДНК.
